\documentclass[11pt,a4paper]{article}

\usepackage[T1]{fontenc}
\usepackage[utf8]{inputenc}
\usepackage{charter}
\usepackage[margin=1.9cm]{geometry}
\usepackage{titlesec}
\usepackage{enumitem}
\usepackage[dvipsnames]{xcolor}
\usepackage{microtype}
\usepackage[hidelinks]{hyperref}

\definecolor{accent}{RGB}{40,72,130}
\hypersetup{colorlinks=true, urlcolor=accent, linkcolor=accent}

\setlength{\parindent}{0pt}
\pagestyle{empty}
\setcounter{secnumdepth}{0}

% Section headings: blue, bold, with a rule underneath.
\titleformat{\section}{\large\bfseries\color{accent}}{}{0pt}{}[{\color{accent}\titlerule[0.8pt]}]
\titlespacing*{\section}{0pt}{11pt}{5pt}

% Two-line entry with a right-aligned column (location / dates).
\newcommand{\entry}[4]{%
  {\bfseries #1}\hfill{#2}\\[1pt]%
  {\itshape #3}\hfill{#4}\par\vspace{2pt}%
}

% One-line entry (title left, links right) with a robust two-column layout
% so a long title never pushes the links onto their own line.
\newcommand{\entrymini}[2]{%
  \noindent\begin{minipage}[t]{0.76\textwidth}\raggedright\bfseries #1\end{minipage}%
  \hfill%
  \begin{minipage}[t]{0.22\textwidth}\raggedleft\small #2\end{minipage}%
  \par\vspace{1pt}}

\newlist{tight}{itemize}{1}
\setlist[tight]{leftmargin=1.35em, itemsep=1.5pt, topsep=2pt, parsep=0pt, label=\textbullet}

\begin{document}

% ---------- Header ----------
\begin{center}
  {\LARGE\bfseries Amitesh Puri}\\[4pt]
  {\small (+91) 7602012472 \quad\textbar\quad
   \href{mailto:amiteshpuri007@gmail.com}{amiteshpuri007@gmail.com} \quad\textbar\quad
   \href{https://www.linkedin.com/in/amitesh-puri-7b6b991aa/}{LinkedIn} \quad\textbar\quad
   \href{https://github.com/AmiteshPuri}{GitHub} \quad\textbar\quad
   \href{https://amiteshpuri.github.io}{Website}}\\[2pt]
  {\small West Bengal, India \textperiodcentered\ 721433}
\end{center}

% ---------- Education ----------
\section{Education}
\entry{Manipal Academy of Higher Education (MIT Manipal)}{Udupi, Karnataka, India}
      {M.Sc. in Applied Mathematics and Computing \quad CGPA: 7.28 / 10}{Aug 2023 -- Jun 2025}
\vspace{4pt}
\entry{Ramnagar College, Vidyasagar University}{Paschim Medinipur, West Bengal, India}
      {B.Sc. in Mathematics \quad Final Grade: 62\%}{Jul 2017 -- Oct 2020}
\vspace{4pt}
\entry{Kishorenagar Sachindra Siksha Sadan}{Contai, West Bengal, India}
      {Higher Secondary (Class XII) \quad Final Grade: 74.5\%}{2014 -- 2016}

% ---------- Research Experience ----------
\section{Research Experience}
\entrymini{Addressing Conditional Shift in Graph Data: An Empirical Study of Distance Metrics}
          {\href{https://amiteshpuri.github.io/thesis/MSc_Project_Report.pdf}{[Report]}\ \
           \href{https://github.com/AmiteshPuri/MSc--Research-Project}{[GitHub]}}
{\itshape Master's Research Report, Manipal Academy of Higher Education \quad Supervisor: Dr. Akansha Singh}\par\vspace{2pt}
\begin{tight}
  \item Studied the robustness of semi-supervised node classification under conditional distributional shift across several GNN architectures on benchmark datasets.
  \item Added distance metrics as auxiliary regularisers that penalise divergence between training and test node representations, isolating each metric's effect in controlled ablations.
  \item Used conformal prediction to evaluate coverage guarantees under distributional change.
  \item Implemented all experiments in PyTorch with fixed random seeds and publicly available code.
\end{tight}

% ---------- Research Projects ----------
\section{Research Projects}
\entrymini{Physics-residual Flow Matching under distribution shift}
          {\href{https://amiteshpuri.github.io/reports/fm-ood-residual.pdf}{[Report]}\ \
           \href{https://github.com/AmiteshPuri/fm-ood-residual}{[GitHub]}}
{\itshape Independent research report}\hfill{2026}\par\vspace{2pt}
\begin{tight}
  \item Measured whether a physics-residual penalty improves a flow-matching generative model, and whether the conflict-free PBFM formulation is worth its added complexity over a plain weighted residual.
  \item Ran the comparison across five seeds, seven sampling budgets, and four graded distribution shifts, and reported the results without selection.
  \item Found that the weighted residual improved every in-distribution metric but did not hold up under shift, where it often reversed.
\end{tight}
\vspace{5pt}
\entrymini{When does a geodesic probability path help latent flow matching?}
          {\href{https://amiteshpuri.github.io/reports/fm-paths.pdf}{[Report]}\ \
           \href{https://github.com/AmiteshPuri/FM-paths}{[GitHub]}}
{\itshape Independent research report}\hfill{2026}\par\vspace{2pt}
\begin{tight}
  \item Tested whether a straight-line or a geodesic probability path changes the outcome of latent flow matching, with the data, model, optimiser, and seed held fixed.
  \item Found the two paths equal across converged, seed-matched runs, a controlled null result.
  \item Showed through an ablation that path geometry helps only when the data lie on the assumed manifold.
\end{tight}

% ---------- Technical Training ----------
\section{Technical Training}
\entry{Guided Projects Student}{Nov 2022 -- Feb 2023}
      {The Machine Learning Company}{Remote}
\begin{tight}
  \item Completed applied projects in NLP and computer vision, including end-to-end pipelines for text classification, sentiment analysis, and language generation.
  \item Trained and evaluated deep learning models for image classification and object detection tasks.
\end{tight}
\vspace{5pt}
\entry{Data Science Trainee}{Aug 2021 -- Jan 2022}
      {Evarcity}{Bangalore, India}
\begin{tight}
  \item Built Python-based data analysis and machine learning workflows on real-world datasets, with practical exposure to feature engineering, model evaluation, and data visualisation.
\end{tight}

% ---------- Personal Projects ----------
\section{Personal Projects}
\entrymini{Cyberbullying Detection}{\href{https://github.com/AmiteshPuri/twitter-sentiment-classification-cyberbullying-}{[GitHub]}}
\begin{tight}
  \item Developed a text classifier in TensorFlow for identifying cyberbullying in social media content. Deployed the model through a Gradio interface, with user feedback collected and managed using Argilla for iterative data annotation.
\end{tight}
\vspace{4pt}
\entrymini{Text Generation with GPT-2}{\href{https://github.com/AmiteshPuri/text-generation-using-gpt2}{[GitHub]}}
\begin{tight}
  \item Fine-tuned GPT-2 on a domain-specific corpus and deployed an interactive text generation interface for user evaluation and qualitative assessment of output quality.
\end{tight}
\vspace{4pt}
\entrymini{English to Hindi Machine Translation}{\href{https://github.com/AmiteshPuri/fine-tuned-translator-eng-hin}{[GitHub]}}
\begin{tight}
  \item Fine-tuned a transformer-based sequence-to-sequence model for English to Hindi translation and deployed the system through a Streamlit application with integrated text-to-speech output.
\end{tight}

% ---------- Technical Skills ----------
\section{Technical Skills}
\begin{tabular}{@{}p{3.9cm}@{\hspace{6pt}}p{12.6cm}@{}}
  \textbf{Programming}        & Python, R \\[2pt]
  \textbf{ML / DL Frameworks} & PyTorch, TensorFlow, Keras, Hugging Face Transformers, Scikit-learn, PyCaret \\[2pt]
  \textbf{Data Analysis}      & Pandas, NumPy, Matplotlib, Seaborn, Data Visualization \\[2pt]
  \textbf{Tools}              & Git, Gradio, Streamlit, Hugging Face Hub, Argilla \\
\end{tabular}

\end{document}
